\section{Secants and terminal curves}\label{geo:section}

We work in the contradiction regime \eqref{red:regime}.
In particular, $N=B^d$, $A=B^2$, and $M=o(A)$.
We use the flat bounds of Lemma~\ref{red:flat-caps}: for some fixed
$c_0>0$ every proper complex affine flat contains at most $NB^{-c_0}$
sample points, a plane contains at most $B^{2+o(1)}$ sample points,
and a line contains at most $2M$ sample points.
The complex-flat assertion follows by taking the real affine span of
its real points. Constants throughout this section depend only on $d$.
For an affine point $p$, write
\[
 \pi_p(y)=[y-p]\in\mathbb P^{d-1},\qquad y\ne p.
\]
All varieties in the geometric arguments are reduced complex varieties;
their degrees are degrees of projective closures.

\subsection{The secant tangent calculation}

The generic conclusion is a form of generalized trisecant geometry;
compare \cite[Section~3.2]{kaminskikanelbelovteicher2008}.  We give the
tangent calculation because the assertion for a particular pair is
also needed.

\begin{lemma}[Secants and a particular-pair version]\label{geo:secants}
Let $X,Y\subset\mathbb A^d$ be irreducible varieties of dimensions
$k,l\ge1$. Suppose that for a general pair $(x,y)\in X\times Y$ the
connecting direction is in neither $T_xX$ nor $T_yY$.
Let $Y'$ be another irreducible $l$-dimensional variety, possibly equal
to $Y$. If a general such line has a further point
$z\in Y'\setminus\{x,y\}$, then $k\le l$.
If $k=l$, then $X$ and $Y$ lie in a common affine $(k+1)$-flat.
More generally, when $k=l$, equality of the projected tangent spaces
at a general pair already implies this common-flat conclusion,
without an assumption about a further root.

There is also the following local assertion. Suppose $p\in X$ and
$q\in Y$ are smooth, and the map
\[
 \psi(x,y)=(x,\pi_x(y))
\]
has differential of rank $k+l$ at $(p,q)$.
Suppose a smooth local variety $W$ of dimension $k+l$ at
$(p,\pi_pq)$ contains both product images
$\psi(X\times Y)$ and $\psi(X\times Y')$ locally at the relevant
pairs. If $z\in Y'\setminus\{p,q\}$ has $\pi_pz=\pi_pq$, then
\[
 \overline{T_pX}\subseteq\overline{T_qY}
 \quad\hbox{in }\mathbb C^d/\mathbb C(q-p).
 \tag{\ref*{geo:secants}.1}\label{geo:tangent-inclusion}
\]
Here bars denote the images of vector tangent spaces. In particular,
if these images have dimensions $k=l$ and are unequal, no such $z$
exists. The same source component $X$ must occur at both roots.
\end{lemma}

\begin{proof}
Choose a direction chart by requiring one coordinate of the direction
to equal one. Write $y-x=a u$ and $z-x=b u$ in that chart; the two
nonzero longitudinal parameters $a,b$ are distinct.
For a source velocity $\xi\in T_xX$, variation of $x$ with the other
endpoint fixed has direction derivative $-\bar\xi/a$ at $y$, and
$-\bar\xi/b$ at $z$. The first coordinates of these two image tangent
vectors are both $\xi$. Their difference is therefore vertical and
has direction component $(b^{-1}-a^{-1})\bar\xi$.

At a general pair the product image has dimension $k+l$: its first
coordinate records the $k$ source directions, and injectivity of the
target tangent modulo the connecting line supplies $l$ vertical
directions. If the further root exists generally, the two irreducible
product images coincide. Indeed they have the same dimension $k+l$,
and the first is contained in the second on a dense set. At a general
smooth image point their tangent spaces agree. The vertical kernel of
the first product image is exactly $\overline{T_yY}$.
The preceding subtraction gives
$\overline{T_xX}\subseteq\overline{T_yY}$ and hence $k\le l$.
This calculation also proves the local assertion: the full-rank first
product identifies its tangent image with $T W$, whose vertical kernel
is $\overline{T_qY}$. Source variation with the second root held fixed
still gives tangent vectors in $T W$. No assertion that the second
root belongs to the finite sample is used.

When $k=l$, the projected tangents agree. Thus
\[
 T_xX\subseteq T_yY+\mathbb C(x-y),
 \qquad
 T_yY\subseteq T_xX+\mathbb C(x-y).
\]
Fix a general $y$. Project $X$ projectively away from the affine tangent
plane $y+T_yY$. The first inclusion says that this projection has zero
differential wherever it is defined. In characteristic zero a rational
map with zero differential has zero-dimensional image: on function
fields, every function in its image has zero differential and is
constant. Its irreducible image is consequently one point.
Thus $X$ lies in an affine $(k+1)$-flat containing $y+T_yY$.
The case in which $X$ already lies in the projection center gives the
same conclusion. Interchanging $X,Y$ bounds both affine spans by
$k+1$. If one span has dimension $k+1$, the preceding flat for each
general point of the other variety must equal that span, so it contains
both varieties. If both spans have dimension $k$, both varieties are
flats, and the displayed tangent inclusions say directly that their
affine join has dimension at most $k+1$.
\end{proof}

\subsection{A generic count for families of carriers}

The next formulation separates the geometric count from the scale
regularization in Section~\ref{sec:profiles}.

\begin{lemma}[Generic projection depth]\label{geo:generic-depth}
Let $\mathcal X,\mathcal Y$ be finite families of distinct irreducible
carriers of dimensions $k\ge l\ge1$, where $k<d$ and $l<d-1$.
Fix assigned sets
\[
 S_X\subseteq P\cap X\quad(X\in\mathcal X),\qquad
 T_Y\subseteq P\cap Y\quad(Y\in\mathcal Y),
\]
pairwise disjoint within each family.  The source and target assignments
may overlap with each other.  If an initial list contains repeated copies
of a carrier, pool those copies and their assigned sets before defining
these distinct families and checking the following hypotheses.  Put
\[
 n_x=\sum_{X\in\mathcal X}|S_X|,\qquad
 n_y=\sum_{Y\in\mathcal Y}|T_Y|.
\]
The normalized component weights are $|S_X|/n_x$ and $|T_Y|/n_y$,
and a component pair has the product of its two weights.
Assume $n_x,n_y\ge N\exp(-o(w))$, with $w\to\infty$ and
$w\le\log B$.  Assume that every flat member is real. Suppose, when
$k=l$, that for some number $e$ the number of flat members in each
family is at most $e\exp(o(w))$, and that the normalized weight of
each such member is at most $e^{-1}\exp(o(w))$.

For every fixed $\epsilon>0$, one can discard component pairs of
total product weight $o(1)$, and obtain for each remaining source
component $X$ a list $\mathcal Y_X$ of target components, such that:
for every $Y\in\mathcal Y_X$, at a general pair $(x,y)\in X\times Y$
both projected tangents are injective and the number of moving roots
on $\bigcup_{Z\in\mathcal Y_X}Z$ over $\pi_xy$ is at most
$\exp(\epsilon w)$. This counts the generic map degrees of all listed
components sharing that projected image, and excludes the undefined
root $x$.

The hypotheses apply to the two carrier configurations used below:
in a common window, $k=l$, $w=s$, and $e=e_j$; in separated subpower
windows, $w=\log B$, and, when $k=l$, $e=B^{d-l}$.
\end{lemma}

\begin{proof}
First discard pairs for which a general connecting line is tangent at
one endpoint. For a fixed target component $Y$, choose a general
smooth $y\in Y$. Every source component tangent in this sense is
contained in the fixed proper affine flat $y+T_yY$.
Its total assigned source mass is at most $NB^{-c_0}$.
After division by $n_x$, this is $B^{-c_0+o(1)}$.
Average with target weights. The argument with source and target
interchanged removes tangencies at the other endpoint. The choices
of tangent points can be made simultaneously for the finite lists.

For $k=l$, also discard pairs lying in a common $(k+1)$-flat when one
member is nonlinear. That nonlinear member spans the flat: its span
cannot have dimension $k$ unless it is itself a flat. Fixing that
member therefore pins down a proper flat containing the other
member's assigned points. The same flat bound and weighted averaging
show that these pairs have product weight $o(1)$.

Lemma~\ref{geo:secants} now rules out a third moving root except when
$k=l$ and the source and first target are flats in a common
$(k+1)$-flat $H$. Every further listed target with the same projection
must also lie in $H$: for a general $x\in X$, its affine join with $x$
equals $H$. It is a flat, since the nonlinear exception was removed.
Each such flat contributes exactly one moving root. Thus it remains
to bound the number of listed flats in $H$ for most weighted pairs.

The flat members are already distinct within each family.  In the slice,
identify a flat occurring in both families and add its two weights; this
changes the bounds by at most a factor two. Fix a small constant
$\eta>0$. Delete flats of normalized weight less than
$e^{-1}\exp(-\eta w)$; their total weight is
$\exp(-\eta w+o(w))$. Slice all remaining flats by a common generic
real projective subspace of codimension $k$. Each becomes a point,
and each generated $(k+1)$-flat becomes a line. For these finite
families, genericity ensures that distinct generated flats yield
distinct lines and that a sliced point belongs to such a line exactly
when the original flat is contained in the corresponding $H$.
The number of sliced points is at most $e\exp(o(w))$.

Let $t_H$ be their number on the line for $H$. The lower weight cutoff
and the physical flat bound give
\[
 t_H\le e\exp(\eta w+o(w))B^{-c_0}.
\]
Choose $\eta<c_0/4$. Since $w\le\log B$, the last bound is at most
$eB^{-c_0/2}$ for large $B$. By the rich-line form of
Szemer\'edi--Trotter, Theorem~\ref{ci:st}, the number of these
lines at dyadic occupancy $t$ is
\[
 O\!\left(\frac{e^2\exp(o(w))}{t^3}
            +\frac{e\exp(o(w))}{t}\right).
\]
Each has pair weight at most $t^2e^{-2}\exp(o(w))$.
Summing for $t\ge\exp(\epsilon w)$ gives
\[
 \exp(-\epsilon w+o(w))+B^{-c_0/2}\exp(o(w))=o(1).
\]
Delete these pairs and, for each source, delete the corresponding
targets from its list. A pair which remains belongs to an $H$ with
at most $\exp(\epsilon w)$ remaining target flats. Removing targets
cannot increase this number. This proves the assertion.

For the announced applications, the degree and per-degree mass bounds
of Proposition~\ref{prof:budget} give the common-window hypotheses.
In subpower windows all testing degrees are $B^{1+o(1)}$, so the
number of $l$-flat carriers is $B^{d-l+o(1)}$ and their normalized
weights are at most $B^{-(d-l)+o(1)}$. These are the second set of
hypotheses. Exceptional intersections of distinct projected images
are not part of this statement: at a generic direction, only
components having the same image contribute.
\end{proof}

\subsection{Deletion of assigned terminal curves}

Fix, for each $p\in P$, a union $C_p$ of irreducible curves in
$\mathbb P^{d-1}$, each spanning a projective subspace of dimension at
most two, with summed degree at most
\[
 E=\delta N/A,\qquad \delta\longrightarrow0.
\]
Identical components are pooled. An assignment below is partial:
some sample points may be left unassigned, and no point is assigned
twice.

\begin{proposition}[Terminal-curve deletion]\label{geo:terminal}
Assign sample points to finitely many irreducible physical curves
$Y_i$, each containing its assigned points. Pool repetitions of a
curve, and require every line member to be real. Then the number of
ordered pairs $(p,q)$, $p\ne q$, for which $q$ is assigned to $Y_i$
and
\[
 \overline{\pi_pY_i}\subseteq C_p
\]
is $o(N^2)$. The estimate is uniform in the assignment and the curves.
\end{proposition}

\begin{proof}
Write $n_i$ for the assigned mass of $Y_i$ and $d_i=\deg Y_i$.
Lemma~\ref{poly:separators} gives
\[
 \sum_i n_i\le N,\qquad n_i\le C M d_i.
 \label{geo:curve-mass}\tag{\ref*{geo:terminal}.1}
\]
We treat lines first, then nonlinear curves.

\paragraph{Assigned lines.}
Centers on their assigned target line contribute at most
$2M\sum_i n_i=o(N^2)$. Otherwise the projected line is itself a
line component of $C_p$, and corresponds to a real plane $H$
containing $p$ and the assigned line. Put
\[
 S_H=\sum_{Y_i\text{ a line},\,Y_i\subset H}n_i,
 \qquad h_H=|P\cap H|.
\]
Planes with $S_H<A$ account for at most $NAE=\delta N^2$ pairs,
because a center has at most $E$ line components in $C_p$.
For the other planes it suffices to prove
$\sum_{S_H\ge A}h_HS_H=o(N^2)$.
They form a finite family determined by pairs of distinct assigned
lines: one line has weight at most $2M=o(A)$.

If an assigned line in $H$ has weight $b$, then
\begin{equation}\label{geo:line-plane-product}
 b h_H\le C M^2.
\end{equation}
To prove this, choose plane coordinates with that line as the first
axis and map a point $(x_1,x_2)$ to $(x_1,x_1^2+x_2^2)$.
Every fiber has size at most two. A distance from an axis point
$(a,0)$ becomes a line equation
$v-2au=r^2-a^2$ in these coordinates. There are at most $bM$ distinct
such lines, with different slopes for different axis points, and the
incidence count is at least $bh_H/2$. The point-line estimate gives
\[
 bh_H\le C\bigl((bh_HM)^{2/3}+h_H+bM\bigr).
\]
For $b$ above a sufficiently large fixed constant, absorb the $h_H$
term. If $h_H=O(M)$, use $b\le2M$; otherwise absorb $bM$ as well
and cube the remaining inequality. Both cases give
\eqref{geo:line-plane-product}. For bounded $b$, the case $h_H\le1$
is immediate. Otherwise choose two distinct
sample points in $H$. Every point of $P\cap H$ is an intersection of
two circles whose radii belong to the at most $M$ distance values;
two distinct-center circles have at most two intersections. Thus
$h_H\le2M^2$, which proves the bounded-$b$ case too.

Group the planes by dyadic values
$h\le h_H<2h$ and $s\le S_H<2s$, where $h,s\gtrsim A$.
For a dyadic line weight $b$, let $N_b$ be the total assigned mass
of all lines with weight in $[b,2b)$. Equation
\eqref{geo:line-plane-product} permits only $b\le L=CM^2/h$.
Let $F$ be the number of planes in this $(h,s)$ group.
A generic hyperplane slice turns the finite line arrangement into
points and the relevant planes into lines, preserving exactly the
needed incidences. A further generic planar projection permits use
of the real point-line estimate. Applying it to each weight class
and multiplying by its weight yields
\[
 Fs\le C F^{2/3}T(L)+CN+CLF,
 \qquad T(L)=\sum_{b\le L}b^{1/3}N_b^{2/3}.
\]
Since $L/s\le CM^2/A^2=o(1)$, the last term is absorbed, giving
\begin{equation}\label{geo:weighted-plane-count}
 F\le C\left(N/s+T(L)^3/s^3\right).
\end{equation}
The contribution of $N/s$ to $\sum hsF$ is at most
$CN B^{2+o(1)}(1+\log N)^2=o(N^2)$, using the plane cap and $d\ge3$.
For the other contribution, H\"older's inequality gives
\[
 T(L)^3\le
 L\left(\sum_{b\le L}(b/L)^{1/2}N_b^2\right)
   \left(\sum_{b\le L}(b/L)^{1/4}\right)^2
 \le C L\sum_{b\le L}(b/L)^{1/2}N_b^2.
\]
Consequently
\[
 \frac{hT(L)^3}{s^2}
 \le\frac{CM^2}{s^2}
       \sum_{b\le L}(b/L)^{1/2}N_b^2.
\]
For each fixed $b$, summing over dyadic $h$, equivalently over inverse
dyadic $L\ge b$, costs a bounded geometric sum. Then
$\sum_bN_b^2\le N^2$, and $\sum_{s\gtrsim A}s^{-2}=O(A^{-2})$.
The result is $O(N^2M^2/A^2)=o(N^2)$, completing the line case.

\paragraph{Nonlinear curves when $d\ge4$.}
If $\overline{\pi_pY_i}$ spans at most a projective plane, the affine
span of $Y_i$ has dimension at most three. A curve with three-dimensional
span requires the center to lie in that fixed three-flat; the flat cap
and \eqref{geo:curve-mass} in the form $\sum n_i\le N$ show that these
pairs cost $o(N^2)$. For a planar curve, centers in its own plane also
cost $o(N^2)$. It remains to consider planar curves viewed from outside
their planes. Projection is then a projective isomorphism of the plane
onto its image and preserves the curve's degree.

At a center $p$, group curves by their common nonlinear projection
$\gamma\subset C_p$, writing $e_\gamma=\deg\gamma$ and
$W_{p,\gamma}=\sum_{i\text{ in the group}}n_i$.
Each individual mass is at most $CMe_\gamma$.
For two different curves in this group, their planes are distinct and
lie, with $p$, in a projective three-space. Their intersection is a
projective line $L$.

For a fixed ordered pair of such curves, there are at most
$C e_\gamma$ possible centers. Here is the required stabilizer argument.
Each perspective between the two planes fixes $L$ pointwise and maps
one curve to the other. Comparing two perspectives gives a projectivity
of the first plane preserving its nonlinear curve and fixing $L$
pointwise. In the affine chart with $L$ at infinity, it has the form
$u\mapsto\lambda u+b$. A nonidentity translation cannot preserve an
irreducible nonlinear curve: the infinite translation orbit of any
point lies on a line, whose closure would be contained in that curve.
Thus the scalar homomorphism on this stabilizer has trivial kernel;
the stabilizer is abelian because every commutator is a translation.
If it is nontrivial, a nonidentity element has a unique affine fixed
point $c$, fixed also by every commuting element. The group consists
of homotheties about $c$. The orbit of a generic point is free and
lies on its line through $c$, so B\'ezout bounds the group order by
the curve degree. This argument also rules out an infinite group.
Finally a perspective between distinct planes determines its center
uniquely: it is the intersection of the joining lines of corresponding
points, two suitably chosen pairs being enough. The claimed bound
on centers follows.

Discard groups with $W_{p,\gamma}\le KMe_\gamma$, for a sufficiently
large fixed $K$. Their total contribution is at most $KNME=o(N^2)$.
In every remaining group the off-diagonal ordered weight satisfies
\[
 \sum_{i\ne j}n_i n_j
   =W_{p,\gamma}^2-\sum_i n_i^2
   \ge\tfrac12 W_{p,\gamma}^2.
\]
If the surviving pair mass were at least $\varepsilon N^2$, weighted
Cauchy--Schwarz, first in $\gamma$ and then in $p$, would give
\[
 \sum_{p,\gamma}\frac{1}{e_\gamma}\sum_{i\ne j}n_i n_j
 \ge \frac{\varepsilon^2N^3}{2E}.
\]
For each fixed ordered pair of distinct curves, degree preservation
and the center bound give
$\sum_p e_\gamma^{-1}\le C$. The left side is therefore at most
$C\sum_{i,j}n_i n_j\le CN^2$, a contradiction since $E=o(N)$.

\paragraph{Nonlinear curves when $d=3$.}
The degree and multiplicity analysis of curves sharing a radial
projection here is closely related to the curve--cone estimates of
Tidor, Yu, and Zakharov
\cite[Sections~10.2--10.3]{tidoryuzakharov2026}.

We first delete centers where
$\operatorname{mult}_pY_i>d_i/2$. There are only boundedly many such
centers for each curve. To see this, take a generic birational plane
projection, preserving or increasing multiplicities at this finite
list of points, as in Lemma~\ref{poly:projection}. For a reduced plane
curve of degree $D$, intersection with a general polar derivative
gives
\begin{equation}\label{geo:polar-bound}
 \sum_{z:\,\operatorname{mult}_z\ge2}
       (\operatorname{mult}_z)^2\le 2D(D-1).
\end{equation}
Indeed the intersection multiplicity at a point of multiplicity $m$
is at least $m(m-1)$, and $m^2\le2m(m-1)$ for $m\ge2$.
A nonlinear irreducible degree-two curve has no point of multiplicity
greater than one; for all other degrees this bounds the number with
multiplicity greater than half the degree by an absolute constant.
The deletion therefore costs $O(\sum n_i)=O(N)$ pairs.
Projection to a line requires the curve and center to lie in one fixed
plane; these pairs cost $o(N^2)$ by the plane cap.

Group the remaining curves at $p$ by their common nonlinear projection
$\gamma$. Let $t_i$ be the generic degree of $\pi_p:Y_i\dashrightarrow
\gamma$ in this group.  The point-projection formula, obtained from the
linear-system degree formula with its base-point contribution
\cite[Proposition~4.4]{fulton1998}, and the multiplicity deletion give
\begin{equation}\label{geo:map-degree-comparison}
 d_i-\operatorname{mult}_pY_i=t_i e_\gamma,
 \qquad d_i\le2t_i e_\gamma.
\end{equation}
An ordered triple $(i,j,k)$ in a group is called valid if each index
occurs at most its corresponding map degree. We claim that for each
fixed ordered triple of curve indices
\begin{equation}\label{geo:triple-center-bound}
 \sum_{p:\,(i,j,k)\text{ valid at }p} e_\gamma^{-2}\le C.
\end{equation}

Choose $v$ general on $Y_i$, simultaneously for the finite list of
centers in this sum. Distinct centers give distinct directions from
$v$, since $v$ can avoid all the lines joining pairs of centers.
The map degrees $t_i,t_j,t_k$ still refer to projection from each
original center $p$; projection from $v$ will be used to count the
directions of those centers.

We first claim that, viewed from $v$, each of $Y_j,Y_k$ projects
birationally to a plane curve, and that the two plane images are
distinct unless $Y_j=Y_k$. A general line joining $v\in Y_i$ to a
point of either target curve is tangent at neither endpoint.
Otherwise, after fixing the tangent endpoint, the other nonlinear
curve would lie in its tangent line. Thus the non-tangency hypothesis
of Lemma~\ref{geo:secants} holds.

Failure of birationality would give a further point on the same
target curve over a general joining direction. Equality of the two
plane images for $Y_j\ne Y_k$ would give a further point on the
other target curve. In either case Lemma~\ref{geo:secants} puts
$Y_i$ and the first target curve in a common plane. At any center
$p$ contributing to the sum, the curves belong to one original
nonlinear group. If the two curves are distinct, this is impossible:
projection from a point in their plane has
a line image, whereas projection from outside that plane is
injective. If the first target is $Y_i$ itself, the conclusion makes
$Y_i$ planar. The repeated first index then makes validity at that
center require
$t_i\ge2$, whereas a planar curve in a nonlinear group is projected
from outside its plane and has map degree one. This proves the claim.

For a valid center $p$, the line $pv$ consequently gives at least
\[
 r_j=t_j-\mathbf1_{i=j},\qquad
 r_k=t_k-\mathbf1_{i=k}
\]
distinct preimages other than $v$ on the respective curves. They can
be chosen generic on those curves by the original choice of $v$.
Birationality identifies the normalizations of the projective closures
of the curve and its plane image. These distinct generic smooth
preimages therefore give distinct normalization points, hence distinct
branches at the corresponding point of the plane image. Validity and
\eqref{geo:map-degree-comparison} imply
\[
 r_j\ge t_j/2\ge d_j/(4e_\gamma),\qquad
 r_k\ge t_k/2\ge d_k/(4e_\gamma).
\]
If the two plane images are distinct, B\'ezout bounds the sum of their
local multiplicity products at these distinct projected centers by
$d_jd_k$. If they agree, then $j=k$, and the residual multiplicity
is at least two: it is $t_j\ge2$ when $i\ne j$, and $t_i-1\ge2$
when all indices agree. Equation~\eqref{geo:polar-bound} bounds the
sum of squared residual multiplicities by $2d_j^2$.
Both cases give \eqref{geo:triple-center-bound}.

Discard once more groups of mass $W_{p,\gamma}\le KMe_\gamma$.
Their contribution is $O(NME)=o(N^2)$.
In a larger group, an invalid ordered triple must repeat an index
whose map degree is at most two. Such an index has mass at most
$4CMe_\gamma$ by \eqref{geo:curve-mass} and
\eqref{geo:map-degree-comparison}. The weight of invalid triples is
at most
\[
 3W_{p,\gamma}\sum_{t_i\le2}n_i^2
 \le 12CMe_\gamma W_{p,\gamma}^2.
\]
Choose $K$ so this is at most half of $W_{p,\gamma}^3$.
If the remaining pair mass were at least $\varepsilon N^2$, H\"older
in each row and then in the centers would imply
\[
 \sum_{p,\gamma}\frac{1}{e_\gamma^2}
       \sum_{(i,j,k)\text{ valid}}n_i n_j n_k
 \ge \frac{\varepsilon^3N^4}{2E^2}.
\]
Indeed $\sum_\gamma W_{p,\gamma}^3/e_\gamma^2\ge
(\sum_\gamma W_{p,\gamma})^3/E^2$, and the sum of the row cubes is
at least their total cubed divided by $N^2$.
On the other hand, \eqref{geo:triple-center-bound} bounds this expression
by $C\sum_{i,j,k}n_i n_j n_k\le CN^3$.
For $d=3$ we have $E=\delta N^{1/3}$, so $N^4/E^2$ exceeds $N^3$
by an unbounded factor. This contradiction proves the proposition.
\end{proof}
